\begin{table}[tp]
\centering
\caption{Stochastic generalized-Lipschitz monotone root finding: \href{https://github.com/k3nfalt/longtails/blob/main/LongTails/Setups/StochasticGeneralizedMonotone.lean}{Stochastic\allowbreak{}Generalized\allowbreak{}Monotone}.}\label{tab:StochasticGeneralizedMonotone}
\parbox{\linewidth}{\footnotesize A continuous monotone map $F:E\to E$ on a finite-dimensional real Hilbert space satisfies $\|F(x)-F(y)\|\le (L_0+L_1\sup_{z\in[x,y]}\|F(z)\|)\|x-y\|$ and has a root $x_\star$ with $\|x_0-x_\star\|\le R$. The jointly measurable oracle uses IID samples, with integrable $g(x,\xi)$ and squared error, $\mathbb{E} g(x,\xi)=F(x)$, and $\mathbb{E}\|g(x,\xi)-F(x)\|^2\le\sigma^2$. Here $L_0,L_1,R,\sigma\ge0$, $\varepsilon>0$, and $0<\delta<1$. Target: $\mathbb{P}(\|F(X)\|>\varepsilon)\le\delta$.}\par\smallskip
\begin{threeparttable}
\footnotesize
\setlength{\tabcolsep}{4pt}\renewcommand{\arraystretch}{1.25}
\begin{tabularx}{\linewidth}{@{}>{\raggedright\arraybackslash}p{0.20\linewidth}>{\raggedright\arraybackslash}p{0.32\linewidth}>{\raggedright\arraybackslash}p{0.25\linewidth}>{\raggedright\arraybackslash}X@{}}
\toprule
\textbf{Result} & \textbf{Oracle complexity} & \textbf{Novelty} & \textbf{Limitations} \\
\midrule
\textbf{(1)} \href{https://github.com/k3nfalt/longtails/blob/main/LongTails/Theorems/StochasticGeneralizedMonotone.lean}{Stochastic\allowbreak{}Generalized\allowbreak{}Monotone} \textbf{(Ours)} $(\star\star\star)$ \href{https://github.com/k3nfalt/longtails/blob/main/scripts/reviews/StochasticGeneralizedMonotone.md}{(review)}
& $\begin{gathered}\mathcal{O}\!\left(\left[1+\tfrac{L_0R}{\varepsilon}+(L_1R)^2+\tfrac{\sigma^2}{\varepsilon^2}\right]\right.\\[-2pt]\left.{}\cdot\ell^3[\ell+\log(1/\delta)]\right)\end{gathered}$\par
All samples, batches and repetitions included.
& ---
& --- \\
\midrule[1.2pt]
Accumulative regularization \citep[Thm.~2.2]{sgmJiLanZhu2026v1}
& $\begin{gathered}\mathcal{O}\!\left(\left[1+\tfrac{L_0R}{\varepsilon}\right](1+\log\ell_0)\right.\\[-2pt]\left.{}+\tfrac{\sigma^2}{\varepsilon^2}\ell_0^3(1+\log\ell_0)^2\right)\end{gathered}$\par
Special case $L_1=0$, $0<\varepsilon\le L_0R$; stochastic calls for $\mathbb{E}\|F(X)\|^2\le\varepsilon^2$.
& i) \textbf{(1)}: scope: permits $L_1>0$.\newline ii) \textbf{(1)}: high probability under finite variance.\newline iii) \textbf{(1)}: matches polynomial accuracy dependence at $L_1=0$.
& \textbf{(1)}: extra logarithmic overhead at $L_1=0$; covers roots only, not constrained VIs. \\
\midrule
Clipped Korpelevich \citep[Thm.~4.4]{sgmVankovNedicSankar2025v2}
& No comparable residual budget: expected squared-distance guarantee, with two stochastic calls per iteration.\par
Restricted to $p$-quasi-sharp operators ($p\ge2$, $\mu>0$) and $0<\alpha\le\tfrac12$ for rates.
& i) \textbf{(1)}: no quasi-sharpness.\newline ii) \textbf{(1)}: residual guarantee at $\alpha=1$; high probability instead of expected distance. Rates not comparable.
& \textbf{(1)}: requires monotonicity; excludes the baseline's constrained and structured non-monotone problems. \\
\midrule
Anchored extragradient \citep{sgmAnchoredRepository2026}
& $\begin{gathered}\mathcal{O}\!\left(1+\tfrac{L_0R}{\varepsilon}\right.\\[-2pt]\left.{}+(1+L_1R)\log\!\left[1+\tfrac{(L_0+L_1G)R}{\varepsilon}\right]\right)\end{gathered}$\par
Exact calls; special case $\sigma=0$, $L_0>0$, $\|F(x_0)\|\le G$.
& i) \textbf{(1)}: scope permits $\sigma>0$.\newline ii) \textbf{(1)}: no initial-residual bound $G$ in stopping budget.
& \textbf{(1)} at $\sigma=0$: quadratic $L_1R$ term and extra logarithms; loses the anytime deterministic last-iterate guarantee; needs $R,\sigma,\varepsilon,\delta$ in updates. \\
\bottomrule
\end{tabularx}
\begin{tablenotes}[flushleft]\footnotesize
\item[] \textbf{Abbr:} $\ell=1+\log(1+L_0R/\varepsilon+L_1R)$; $\ell_0=1+\log(1+L_0R/\varepsilon)$; $G$ bounds $\|F(x_0)\|$; $[x,y]$ is the line segment. In the quasi-sharp baseline, $X^\star$ is the solution set, $p$ and $\mu$ satisfy $\langle F(x),x-x_\star\rangle\ge\mu\,\operatorname{dist}(x,X^\star)^p$ for all $x_\star\in X^\star$, and $\alpha$ is the exponent on $\|F\|$ in generalized Lipschitz continuity. Costs count individual oracle evaluations; $\log$ is natural.
\end{tablenotes}
\end{threeparttable}
\end{table}
\clearpage
